% Replacement root-transport passage for Section 5.

Put
\[
  \Phi_n(k):=L_{S_4}(n,k)+(\log 2)k.
\]
Then $\Phi_n(r_4^{\mathrm{co}})=0$ by definition. To evaluate the same
function at the unrestricted root, let
\[
  T_+(n):=\alpha-\frac{n}{r_+(n)}
\]
and introduce the finite-$n$ support loss
\[
  D_{4,n}(T):=
  \mathcal F_{n,S_+}(T)-\mathcal F_{n,S_4}(T).
\]
Since $L_{S_+}(n,r_+)=0$, the exact support-comparison identity (3.8) gives
\begin{equation*}
\begin{aligned}
  \Phi_n(r_+)
  &=L_{S_4}(n,r_+)-L_{S_+}(n,r_+)+(\log 2)r_+\\
  &=r_+\{\log 2-D_{4,n}(T_+(n))\}.
\end{aligned}
\tag{5.9a}
\end{equation*}
This equality is finite and contains no limiting substitution.

We now transport its target mean to the phase center. Equation (3.6) gives
\begin{equation*}
  T_+(n)=T_0+O\!\left(\frac{\log\log n}{\log n}\right)
  \tag{5.9b}
\end{equation*}
uniformly in the phase. On the compact interval $K_*$, equation (3.9) gives
uniform convergence of both finite dual values. The limiting dual value is
differentiable, with
\[
  \mathcal F_S'(T)=-\lambda_S(T),
\]
because the envelope theorem cancels the derivative of the optimizing tilt.
The tilts are uniformly bounded on $K_*$, so both $\mathcal F_{S_+}$ and
$\mathcal F_{S_4}$ are uniformly Lipschitz there.

More explicitly, choose a deterministic sequence
\[
\begin{split}
  \omega_n^{\mathrm{root}}:=
  &\sup_{T\in K_*}
   |\mathcal F_{n,S_+}(T)-\mathcal F_{S_+}(T)|\\
  &+\sup_{T\in K_*}
   |\mathcal F_{n,S_4}(T)-\mathcal F_{S_4}(T)|
   +C|T_+(n)-T_0|,
\end{split}
\]
where $C$ is a common Lipschitz constant for the two limiting dual values.
Then $\omega_n^{\mathrm{root}}\to0$ uniformly in the phase, and
\[
  |D_{4,n}(T_+(n))-D_4(\delta)|
  \le \omega_n^{\mathrm{root}}.
\]
Substitution into (5.9a) yields the theorem-facing estimate
\begin{equation*}
  L_{S_4}(n,r_+)+(\log 2)r_+
  =r_+\{\log 2-D_4(\delta)+O(\omega_n^{\mathrm{root}})\}
  =r_+\{\log 2-D_4(\delta)+o(1)\}.
  \tag{5.10}
\end{equation*}
Thus the $o(1)$ in (5.10) is one deterministic error sequence valid across
the complete phase, including integer sequences approaching either endpoint.

Finally, Lemma~5.1 gives
$\log 2-D_4(\delta)\ge\gamma_4>0$. Hence (5.10) implies
$\Phi_n(r_+)>0$ for all sufficiently large $n$, uniformly in the phase. The
derivative estimate (3.7) is positive throughout the common root corridor, so
$\Phi_n$ is strictly increasing there. Since
$\Phi_n(r_4^{\mathrm{co}})=0$, this proves
\[
  r_4^{\mathrm{co}}<r_+.
\]
The interval used in the subsequent mean-value argument is therefore
nonempty, with its orientation fixed.
